% TOP_PROBLEM: {"id":30007543,"problem_number":"LOCAL-30007543","title":"Misallocation versus measurement error","category_id":21,"category":{"id":21,"name":"economics","display_name":"Economics","description":"Open economic theory statements, published research agendas, and proposed scoped specifications; question provenance and historical priority are qualified per record.","slug":"economics","order_index":21,"created_at":"2026-10-08T00:00:00Z"},"status":"provenance_pending","external_url":null,"legacy_ids":[],"published":false,"statement":"Identify counterfactual productivity gains from removing specified plant-level wedges separately from measurement error and unavoidable costs.","economics":{"original_id":32,"field":"Growth and productivity","kind":"identification","formalization_basis":"proposed_specification","source_status":"not_verified","priority_status":"not_established","priority_note":"No explicit primary open-problem statement matching this catalogue item was verified; supporting research is not evidence of first formulation.","earliest_verified_reference":null,"current_reference":{"authors":["Mark Bils","Peter J. Klenow","Cian Ruane"],"title":"Misallocation or Mismeasurement?","year":2020,"url":"https://www.census.gov/library/working-papers/2020/adrm/CES-WP-20-07.html","doi":"10.3386/w26711","locator":"Abstract, Census CES-20-07, February 2020","evidence_summary":"The authors propose a measurement-error correction and quantify how it changes reallocation gains. The abstract establishes a specific method and results, not an explicitly open general decomposition problem."},"formalization_note":"New identification/bounding target. The cited paper proposes a correction and obtains results; no exact general open statement was verified."},"statusQualification":"Provenance pending: an explicit published statement of this exact open question was not verified. This is a catalogue-authored candidate specification, not a certified open conjecture. New identification/bounding target. The cited paper proposes a correction and obtains results; no exact general open statement was verified.","statusReviewedAt":"2026-10-08"}
% FIRST_POSTED: 2026-10-08
% ENTRY_KIND: statement_only
% !TeX program = lualatex
\documentclass[11pt]{article}
\usepackage{fontspec}
\usepackage[a4paper,margin=25mm]{geometry}
\usepackage{amsmath,amssymb,mathtools}
\usepackage{xurl}
\usepackage[hidelinks]{hyperref}
\newcommand{\sref}[2]{\hyperlink{source:#1:#2}{[S#2]}}
\setlength{\emergencystretch}{3em}
\begin{document}
% problemId:30007543
\section*{Misallocation versus measurement error}\label{30007543}
\subsection{Definitions and mathematical statement}
Consider plants \(i\) producing a differentiated product with technology \(q_i=A_i k_i^\alpha \ell_i^\gamma m_i^{1-\alpha-\gamma}\), where capital, labor, and intermediates are positive, \(\alpha,\gamma>0\), \(\alpha+\gamma<1\), and production elasticities are specified or identified. Product demand has a stated elasticity, and observed revenue and inputs equal true values plus measurement errors in logs. Economic wedges may affect financing, labor, or intermediate purchases; adjustment costs and transport costs are modeled separately rather than automatically called distortions.
Require initial aggregate output \(Y>0\). Let \(Y^{\mathrm{eff}}\) be aggregate output after eliminating the specified wedges while holding true technologies, aggregate resources, and unavoidable costs fixed, and let \(Y\) be initial output. Define
\[
G=\frac{Y^{\mathrm{eff}}-Y}{Y}.
\]
Identify or bound \(G\) using repeated plant observations, independent audits, or instruments that distinguish measurement error from genuine marginal-product gaps. State error persistence, correlation, and demand assumptions explicitly; without them, dispersed measured revenue products need not identify misallocation. Quantify how much of observed dispersion comes from true wedges, measurement, demand heterogeneity, and adjustment costs. Validate the decomposition using independent responses to resource reallocations. The target is a defensible counterfactual productivity gain in a declared industry, not a presumption that all within-industry productivity dispersion indicates inefficient allocation.

\par\textbf{Formulation and scope.} New identification/bounding target. The cited paper proposes a correction and obtains results; no exact general open statement was verified.

\par\textbf{Open-question provenance.} An explicit published open-question statement was not verified. Sources below are supporting literature and must not be treated as a first listing.

\par\textbf{Historical priority.} No explicit primary open-problem statement matching this catalogue item was verified; supporting research is not evidence of first formulation.

\subsection{Short English statement}
Identify counterfactual productivity gains from removing specified plant-level wedges separately from measurement error and unavoidable costs.

\subsection{Sources}
\begin{itemize}\raggedright
\item[\textnormal{[S1]}]\hypertarget{source:30007543:1}{} Mark Bils, Peter J. Klenow, Cian Ruane. 2020. Misallocation or Mismeasurement?. Abstract, Census CES-20-07, February 2020.
\url{https://www.census.gov/library/working-papers/2020/adrm/CES-WP-20-07.html}.
DOI: 10.3386/w26711.
{\small\textit{Supporting or current literature; see evidence qualification. The authors propose a measurement-error correction and quantify how it changes reallocation gains. The abstract establishes a specific method and results, not an explicitly open general decomposition problem.}}
\end{itemize}
\end{document}
